\documentclass[11pt]{article}
\usepackage[T1]{fontenc}
\usepackage{lmodern}
\usepackage[margin=1.05in]{geometry}
\usepackage{microtype}
\usepackage{amsmath,amssymb,amsthm,mathtools,booktabs}
\usepackage[hidelinks]{hyperref}
\hypersetup{pdftitle={Maximal order of Hardy--Dirichlet coefficient norms and a sharp Hardy coefficient limit}}
\newtheorem{theorem}{Theorem}
\newtheorem{proposition}[theorem]{Proposition}
\newtheorem{lemma}[theorem]{Lemma}
\newtheorem{corollary}[theorem]{Corollary}
\title{Maximal order of Hardy--Dirichlet coefficient norms\\and a sharp Hardy coefficient limit}
\author{}
\date{28 November 2026}
\begin{document}
\maketitle
\begin{abstract}
We prove the square-free maximal-order conjecture of
Bondarenko, Brevig, Saksman, and Seip for coefficients in
\(\mathcal H^p\) of Dirichlet series, throughout \(0<p<1\).
We give a two-term maximal-order expansion, whose second term
is larger in the unrestricted case. The essential disk comparison is
\(C(k,p)\le C(1,p)^k\), with equality exactly for \(k=1,2\).
The proof combines two low-degree eliminants, a sharp entropy
bound, and Kulikov's contractive Hardy-to-Bergman embedding.
We also identify the leading constant in the high-degree
disk coefficient estimate through an exact Fourier extremum
and give an explicit relative \(O_p(n^{-1/2})\) bound.
\end{abstract}

\section{Notation and the two Hardy--Littlewood constants}

For the disk Hardy quasi-norm, define
\[
 C(k,p)=\sup_{\|f\|_{H^p}\le1}|[z^k]f|,
 \qquad \tau=1/p-1,\qquad 0<p<1 .
\]
The general determinant formula in \cite{General} evaluates these
constants by a finite root procedure. The simpler cubic and
sextic formulas in \cite{Low} will be used below.

The sharp constant multiplying \(n^\tau\) in a uniform
coefficient bound is the classical \(C(1,p)\):
\[
 \sup_{n\ge1}\frac{C(n,p)}{n^\tau}
   =C(1,p)=\sqrt{\frac2p}\left(1-\frac p2\right)^{1/p-1/2}.
\]
Indeed for \(\omega^m=1\) the average of \(f(\omega z)\)
keeps precisely the multiples of \(m\) in its Taylor series.
The subadditivity of the \(p\)th power gives its quasi-norm at
most \(m^\tau\|f\|_{H^p}\). Writing the result as a function
of \(z^m\) proves \(C(mk,p)\le m^\tau C(k,p)\).
The displayed sharp constant follows by taking \(k=1\);
the first coefficient was computed in \cite{BBSS,BrevigSaksman}.
Writing \(r=1/p\), the logarithm of that constant has derivative
\(1/r+\log(1-1/(2r))>0\) for \(r>1\).
For example, with \(u=1/(2r)<1/2\) use
\(-\log(1-u)<u/(1-u)<2u\).
The constant equals \(1\) when \(p=1\), so it exceeds
\(1\) throughout \(0<p<1\).

The \(C_p=1\) conjecture discussed in \cite[Section~5]{BrevigSaksman}
is the contractivity of \(H^p\) into \(A^1_{1/p}\).
Kulikov proved this with the normalized weighted area measure
\cite[Corollary~1.3]{Kulikov}. For later use this gives the
coefficient estimate
\begin{equation}\label{eq:gamma-bound}
 C(k,p)\le B_k(r)
 :=\frac{\Gamma(r+k/2)}{\Gamma(r)\Gamma(1+k/2)},
 \qquad r=1/p .
\end{equation}
Indeed for each radius \(t\) the angular \(L^1\) norm of
\(f(te^{i\theta})\) is at least \(|[z^k]f|t^k\).
The \(k\)th radial moment of the normalized measure
\((r-1)(1-|z|^2)^{r-2}\,dx\,dy/\pi\) is \(1/B_k(r)\).
The constant \(1\) for the embedding is attained by constant
functions.

\section{The sharp high-degree coefficient asymptotic}

Fix \(0<p<1\), write \(\tau=1/p-1\), and put
\(c_n=C(n,p)/n^\tau\) for \(n\geq1\). Equip the real line with
\(dx/(2\pi)\) and use the Fourier transform
\(\widehat F(\xi)=(2\pi)^{-1}\int_{\mathbb R}F(x)e^{-i\xi x}\,dx\).
Let \(\mathcal A_p\) be the functions in \(L^1(\mathbb R)\cap
L^p(\mathbb R)\) whose Fourier transforms have compact support
in \([0,\infty)\). Define the half-line Fourier coefficient norm
\[
 A_p=\sup_{\substack{F\in\mathcal A_p\\F\ne0}}
       \frac{|\widehat F(1)|}{\|F\|_{L^p(dx/(2\pi))}}.
\]

\begin{theorem}\label{thm:hl-asymptotic}
The constant \(A_p\) is positive and finite, and
\[
 A_p=\inf_{n\geq1}c_n=\lim_{n\to\infty}c_n,
 \qquad
 C(n,p)=A_p n^\tau\bigl(1+O_p(n^{-1/2})\bigr).
\]
More explicitly, take any integer \(m\) with \(p(2m-1)>1\) and put
\[
 D_m=\frac{6(2m)^{2m}}{11(2m-1)},\qquad
 B_{m,p}=1+\frac{D_m^p}{p(2m-1)-1}.
\]
For every \(n\geq1\) and every \(\epsilon>0\),
\begin{equation}\label{eq:hl-effective}
 \frac{C(n,p)}
 {n^\tau(1+\epsilon)^\tau
      (1+B_{m,p}/(\pi\epsilon n))^{1/p}}
 \ \leq A_p\ \leq\frac{C(n,p)}{n^\tau}.
\end{equation}
In particular, at every algebraic \(p\in(0,1)\) the sharp limiting
constant is computable to any prescribed precision using the exact
finite-degree coefficient calculation.
\end{theorem}

\begin{proof}
If \(F\in\mathcal A_p\), periodize it with period \(L=2\pi n\).
The sum \(P(x)=\sum_{k\in\mathbb Z}F(x+kL)\) converges absolutely
almost everywhere and in \(L^1\) on a period. Its \(j\)th Fourier
coefficient at frequency \(j/n\) is \(n^{-1}\widehat F(j/n)\).
Only finitely many can be nonzero and all have \(j\geq0\).
Fourier uniqueness consequently represents \(P(x)=f(e^{ix/n})\)
almost everywhere for an analytic polynomial \(f\). Subadditivity
of the \(p\)th power gives
\[
 \|f\|_{H^p}^p=\frac1L\int_0^L |P(x)|^p\,dx
 \leq\frac1L\int_{\mathbb R}|F(x)|^p\,dx
 =\frac1n\|F\|_{L^p(dx/(2\pi))}^p.
\]
Since \([z^n]f=n^{-1}\widehat F(1)\), this proves \(A_p\leq c_n\).
A smooth Fourier transform supported near \(1\) and nonzero at
\(1\) furnishes a Schwartz test function, hence \(A_p>0\).

For the reverse bound, let
\[
 K(t)=\kappa_m
  \left(\frac{\sin(t/(2m))}{t/(2m)}\right)^{2m},
 \qquad \int_{\mathbb R}K(t)\,dt=1 .
\]
Its Fourier transform is supported in \([-1,1]\), by convolution
of the Fourier transforms of the sinc factors. For \(|t|\leq1\)
the unnormalized integrand is at least \(11/12\), by
\(\sin u/u\geq1-u^2/6\) there and Bernoulli's inequality.
Thus \(\kappa_m\leq6/11\), and for \(d\geq1\)
\(\int_d^\infty K(t)\,dt\leq D_m d^{-(2m-1)}\).
Put \(K_\epsilon(x)=\epsilon K(\epsilon x)\) and let
\(W=\mathbf1_{[0,L]}*K_\epsilon\), using ordinary convolution
with respect to \(dx\). Positivity and the unit integral of
\(K_\epsilon\) give
\[
   0\leq W\leq1,\qquad
   \sum_{k\in\mathbb Z}W(x+kL)=1 .
\]
For \(x\) outside \([0,L]\), \(W(x)\) is bounded by the one-sided
tail of \(K_\epsilon\) at distance \(\operatorname{dist}(x,[0,L])\).
Consequently, the integral over \([0,L]\) is at most \(L\) and
the two tails give
\begin{equation}\label{eq:hl-window-mass}
 \int_{\mathbb R} W(x)^p\,dx\leq L+\frac{2B_{m,p}}{\epsilon}.
\end{equation}

Let \(f\) be an analytic polynomial and \(P(x)=f(e^{ix/n})\).
For every shift \(h\), the periodized window being identically
one implies
\[
 \frac1{2\pi}\int_{\mathbb R}P(x)W(x-h)e^{-ix}\,dx
       =n[z^n]f.
\]
The average over \(0\leq h\leq L\) of
\(\|PW(\cdot-h)\|_{L^p(dx/(2\pi))}^p\) is
\((2\pi)^{-1}\|f\|_{H^p}^p\int_{\mathbb R}W^p\).
One shift therefore has this \(p\)-power at most
\(n(1+B_{m,p}/(\pi\epsilon n))\|f\|_{H^p}^p\)
by \eqref{eq:hl-window-mass}.
The function \(F_h(x)=e^{i\epsilon x}P(x)W(x-h)\) belongs
to \(\mathcal A_p\): \(W\) is integrable and has bandwidth
\(\epsilon\), and the exponential shifts the entire Fourier
support into \([0,\infty)\).
The sharp norm of evaluation \(\widehat F(\xi)\) on
\(\mathcal A_p\) for \(\xi>0\) is \(\xi^\tau A_p\), by dilation.
Applying this with \(\xi=1+\epsilon\) to \(F_h\) proves the left
inequality of \eqref{eq:hl-effective} for analytic polynomials,
and their density in \(H^p\) yields the assertion for \(C(n,p)\).
Set \(\epsilon=n^{-1/2}\) to obtain the asymptotic and the
computability assertion.
\end{proof}

For completeness there is also an elementary convergence proof.
Averaging \(f\) over rotations by the \(k\)th roots of unity and
rescaling the argument gives
\(C(kn,p)\leq k^\tau C(n,p)\). Also \(C(n,p)\) is nondecreasing
in \(n\) by multiplication by \(z\). For a fixed index \(r\),
compare \(n\) with \(r\lceil n/r\rceil\) to obtain
\(\limsup_n c_n\leq c_r\), which proves convergence to
\(\inf_r c_r\).

The normalized asymptotic constant satisfies the explicit bounds
\[
 \frac{(2/(ep))^{1/p}}{\Gamma(2/p)}
 \ \leq A_p\ \leq\min\left\{\frac{2^{1-1/p}}{\Gamma(1/p)},\ (2/\pi)^\tau\right\}.
\]
For the lower bound, the \(H^p\)-unit test function
\((1-r^2)^{1/p}(1-rz)^{-2/p}\) gives
\[
 C(n,p)\geq \frac{(2/p)_n}{n!}
  \left(\frac{n}{n+2/p}\right)^{n/2}
  \left(\frac{2/p}{n+2/p}\right)^{1/p},
\]
where the choice \(r^2=n/(n+2/p)\) maximizes its \(n\)th
coefficient in this family. Stirling's formula gives the lower
bound. The upper bound follows from the contractive
\(H^p\)-to-\(A^1_{1/p}\) inclusion and radial Fourier extraction.
Indeed the radial moment of order \(n\) of the normalized
weighted Bergman measure is
\(\Gamma(1/p)\Gamma(n/2+1)/\Gamma(n/2+1/p)\).

The second upper bound above follows by taking \(n\to\infty\)
in \(C(n,p)<(\csc(\pi/(2(n+1))))^\tau\) from \cite{General}.
That same bound gives \(C(n,p)<n^\tau\) for every \(n\ge2\):
concavity of sine on \([0,\pi/6]\) gives
\(\sin(\pi/(2(n+1)))\ge3/[2(n+1)]\ge1/n\).

\section{Exact coefficient norms for Dirichlet series}
\label{sec:dirichlet}

Write \(D(N,p)\) for the norm of the \(N\)th coefficient
functional on the ordinary Hardy space \(\mathcal H^p\) of
Dirichlet series, for \(0<p<1\). The Bohr lift identifies a
Dirichlet polynomial with an analytic polynomial on the product
of circles. If \(q_j\) are the successive primes and
\(N=\prod_j q_j^{\kappa_j}\), the \(N\)th coefficient becomes
the coefficient of \(\prod_j z_j^{\kappa_j}\). We use the
Hardy quasi-norm defined by the closure of the analytic
polynomials in boundary \(L^p\); the equivalence with the
definition by finite sections is established in
\cite[Section~2]{BBSS}.

\begin{proposition}[Bondarenko--Brevig--Saksman--Seip]
\label{prop:dirichlet-product}
For every \(N\geq1\) and \(0<p<1\),
\[
 D(N,p)=\prod_{q^k\parallel N} C(k,p).
\]
In particular the one-variable exact formulas determine the
individual coefficient norm at every ordinary integer \(N\).
The norm is attained by the Bohr lift of a product of
one-variable extremizers.
\end{proposition}

\begin{proof}
The product identity is proved in \cite[Section~4]{BBSS}.
For completeness, in an analytic polynomial \(F\) on a finite
polytorus, first extract the \(\kappa_j\)th coefficient in one
variable for each fixed value of the other variables.
By the disk estimate, its absolute value to the power \(p\)
is at most \(C(\kappa_j,p)^p\) times the integral of
\(|F|^p\) in that variable. Integrating over the remaining
variables proves the corresponding operator estimate on
polynomials. Iteration gives the product upper bound, with
\(C(0,p)=1\) by subharmonicity of \(|F|^p\). It is independent
of the number of variables used to represent the polynomial,
so it extends by density. For the reverse bound take the
product of the one-variable norm-one extremizers in the
coordinates with \(\kappa_j>0\), and use Fubini to compute both
its norm and its indicated coefficient. Each factor can be approximated by analytic polynomials
in \(H^p\), so this
finite product belongs to the closure defining the space.
\end{proof}

\begin{corollary}
For rational \(p\in(0,1)\), every \(D(N,p)\) is algebraic.
For algebraic irrational \(p\in(0,1)\), \(D(N,p)\) is
transcendental whenever \(N>1\).
\end{corollary}

\begin{proof}
In the one-variable algebraic candidate formula, write a
winning value as \(d^{-1}U^{-(1/p-1)}\), with \(d>0\), \(U>1\)
both algebraic. A finite product has the same form, taking
the products of the \(d\)'s and of the \(U\)'s. For \(N>1\)
there is at least one factor, so the latter product exceeds
one. The assertions follow by algebraicity of rational powers
and the Gelfond--Schneider theorem, respectively.
\end{proof}

\section{The maximal-order coefficient conjecture}
\label{sec:maximal-dirichlet}

Define \(R_p=\sup_{k\geq1}k^{-1}\log C(k,p)\).
The classical Hardy--Littlewood estimate
\(C(k,p)\ll_p k^{1/p-1}\), together with \(C(1,p)>1\), makes
this a maximum over finitely many \(k\). The conjecture after
equation (34) of \cite[Section~4]{BBSS} asks whether the maximal
order of \(D(N,p)\) over all \(N\) agrees with its maximal order
over square-free \(N\). The latter is \(\log C(1,p)\) in
the following normalization.

\begin{proposition}\label{prop:extremal-prime-power}
For every \(0<p<1\),
\[
 \limsup_{N\to\infty}
       \frac{\log D(N,p)}{\log N/\log\log N}
   =\max_{k\geq1}\frac{\log C(k,p)}{k}.
\]
Consequently the conjecture is equivalent to
\(C(k,p)\leq C(1,p)^k\) for every \(k\geq1\). At each fixed
algebraic \(p\) this condition is decidable by the
one-variable exact algorithm and a finite tail bound.
\end{proposition}

\begin{proof}
For fixed \(k\), take \(N\) to be the \(k\)th power of the
product of the first \(m\) primes. The prime number theorem
gives \(\log N\sim km\log m\), and
Proposition~\ref{prop:dirichlet-product} gives
\(\log D(N,p)=m\log C(k,p)\), yielding the lower bound.

For the upper bound put \(L=\log N\) and
\(y=L/(\log L)^3\). Each exponent of \(N\) is at most
\(L/\log 2\). The total contribution of prime factors
\(q\leq y\) is therefore \(O_p(y\log L)=o(L/\log L)\)
by the classical coefficient estimate. For \(q>y\), the
contribution \(k^{-1}\log C(k,p)\) per unit of exponent
is bounded by \(R_p\), while \(k\log q\geq k\log y\).
The sum of those contributions is at most \(R_pL/\log y\).
Since \(\log y\sim\log L\), this proves the asserted formula.
The same argument restricted to square-free integers uses
only \(k=1\), so the equality condition follows.

For an explicit tail bound, use
\(C(k,p)\leq C(1,p)k^{1/p-1}\), which follows from Wiener's
averaging trick or \cite[Section~5]{BrevigSaksman}. Since
\(\log x/(x-1)\) decreases to zero for \(x>1\), one can find
an integer \(K\) such that
\((1/p-1)\log K/(K-1)<\log C(1,p)\). This strict condition
continues to hold for \(k\geq K\), so only \(k<K\) have to be
compared. For rational \(p\) the comparisons are between
algebraic numbers. If \(p\) is algebraic irrational, write
the candidates in the form \(d^{-1}U^{-(1/p-1)}\) with
\(d,U\) algebraic. Equality of two values implies equality
of the respective \(d\)'s and \(U\)'s by
Gelfond--Schneider, and an unequal pair can be compared by
certified approximation.
\end{proof}

\subsection{Degree three and the square-free comparison}
\label{sec:cubic-domination}

The exact degree-two formula of \cite[Theorem~1]{BrevigSaksman}
gives \(C(2,p)=C(1,p)^2\). In degree three the inequality is
strict.

\begin{lemma}\label{lem:cubic-domination}
For every \(0<p<1\) one has \(C(3,p)<C(1,p)^3\).
\end{lemma}

\begin{proof}
Use the exact cubic of \cite{Low} with \(a=2/p-1\). Its
distinguished root \(u\) is the unique root in \((0,2a)\) of
\[
 P_a(u)=\frac{a-2}{6}u^3+\frac{a(5-a)}4u^2
                  +\frac{a(1-2a)}2u-a^2 .
\]
It satisfies \(u>10a/7\): indeed,
\[
 P_a(10a/7)=-a^2(25a^2-155a+294)/1029<0
\]
by the negative discriminant of the displayed quadratic,
and \(P_a(2a)>0\). Set \(v=u/(2a)\) and \(w=1+(a-1)v\).
Then \(5/7<v<1\), \(w<a\), and the formula yields
\[
 C(3,p)^2=a^2 w U^{1-a},\qquad
 U=1+\frac{1}{a^2w}+\frac{4v^2(w+1)}{w^2}.
\]
For the first coefficient, \(C(1,p)^2=a(1+1/a)^{1-a}\).
Write \(S=(1+1/a)^3\). It is enough to show
\(U/S>(w/a)^{1/(a-1)}\).
The inequality \(\log(1-t)\leq-t\) for \(0<t<1\), followed
by the alternating exponential estimate, gives
\[
 (w/a)^{1/(a-1)}
 \leq e^{-(1-v)/a}
 \leq T:=1-\frac{1-v}{a}+\frac{(1-v)^2}{2a^2}.
\]
To check \(U/S>T\), put
\(J=2a^2(a+1)^3 w^2(U/S-T)\). This is a polynomial.
Substitute \(a=1+B\) and \(v=(5+7t)/(7(1+t))\).
The coefficients of \(7^4(1+t)^4J\) in rows \(B^0,\ldots,B^6\)
and columns \(t^0,\ldots,t^4\) are
\[
 \begin{pmatrix}
 196&8624&26264&27440&9604\\
 2898&51996&148568&152292&52822\\
 10470&127316&343784&347116&120050\\
 16440&159712&413364&414344&144060\\
 12336&106120&269304&271656&96040\\
 4030&34132&88396&91924&33614\\
 350&3780&10976&12348&4802
 \end{pmatrix}.
\]
All entries are positive. Because \(a>1\) and \(5/7<v<1\),
the substitution has \(B,t>0\) and proves the result.
\end{proof}

\subsection{Degree four}

\begin{lemma}\label{lem:quartic-domination}
For every \(0<p<1\), \(C(4,p)<C(1,p)^4\).
\end{lemma}

\begin{proof}
Let \(a=2/p-1>1\) and \(\delta=a-1\), and use the
distinguished parameter \(v\) in the fourth coefficient
theorem of \cite{Low}. Put
\[
\begin{gathered}
 k_0=\frac{\delta v}{2(1-v)},\qquad
 u=\frac{6(1-v)^2(1+v)}{v\delta((a-2)+(2a-1)v)},\\
 S=1+u+(k_0u)^2+v^2u+(v/a)^2,\qquad S_0=(1+1/a)^4 .
 \end{gathered}
\]
The value formula is \(C(4,p)=aS^{-\delta/2}/v\).
Thus with \(Y=S/S_0\) we have
\[
 \log\frac{C(1,p)^4}{C(4,p)}
  =\log(av)+\frac{\delta}{2}\log Y
  \ \geq\ \frac{av-1}{av}+\frac{\delta}{2}\frac{Y-1}{Y},
\]
using \(\log x\ge1-1/x\). The last expression is positive
exactly when
\[
 ((a+1)av-2)S-\delta av S_0>0.
\]
Let \(G_a(v)\) be the numerator of this rational expression,
with denominator \(a^3v\delta((a-2)+(2a-1)v)^2\).
This denominator is positive at the distinguished root.
Also put \(A_j=\binom aj\) and recall the polynomial
\[
 N_a(v)=2av\delta((a-2)+(2a-1)v)^2
  \left(\frac va\left[v+2A_2vu+
           (A_2k_0^2+3A_3k_0+A_4)u^2\right]-1\right).
\]
The distinguished parameter satisfies \(N_a(v)=0\).
Exact elimination gives
\[
 \operatorname{Res}_v(N_a,G_a)
 =59049 a^7(a-2)(a-1)^{18}(a+1)^{12}(2a-1)\,T(a),
\]
where \(T\) has degree twenty-three and all coefficients of
\(T(1+t)\) are positive. At \(a=2\),
\(\gcd(N_2,G_2)\) is a nonzero multiple of \(v\).
Also \(N_2\) changes sign on
\((60969/10^5,60970/10^5)\), the interval lies within the
distinguished branch's root interval, and \(G_2\) is
strictly positive there by monomial interval bounds.
The appendix verifies these identities and signs using
exact polynomial arithmetic.

The parameter \(v\) is continuous for \(a>1\).
For completeness, the normalized extremal factors belong
to the compact family \(Q=P/\|P\|_2\) with
\(P=\prod_{j=1}^4(1-\lambda_jz)\), \(|\lambda_j|\le1\).
On this family the Taylor coefficient modulus of \(Q^{a+1}\)
is jointly continuous in the parameter and coefficients,
by the finite Taylor expansion at zero, with \(Q(0)\)
uniformly bounded away from zero. All limit points of
optimizing factors are therefore optimizing at the limit
parameter. The uniqueness up to rotation in \cite{Low}
identifies the same modulus of the leading coefficient for
every limit, hence the same \(v\). The distinguished
\(G_a(v)\) is consequently continuous. By the resultant
and gcd it never vanishes, and its value is positive at
\(a=2\). Thus the logarithmic lower bound is positive
throughout \(a>1\).
\end{proof}

\subsection{Degree five and the uniform tail}

\begin{lemma}\label{lem:tail-domination}
For every \(0<p<1\) and \(k\ge5\), \(C(k,p)<C(1,p)^k\).
\end{lemma}

\begin{proof}
Put \(r=1/p>1\), \(u=1/(2r)\) and
\(g(r)=\log C(1,1/r)=\tfrac12\log(2r)+(r-\tfrac12)\log(1-u)\).
We use the degree-uniform entropy bound from \cite{General}
in degree five:
\begin{equation}\label{eq:csc5}
 C(5,1/r)<\bigl(\csc(\pi/12)\bigr)^{r-1}.
\end{equation}
We have \(g(1)=0\) and \(g''(r)=(1-r)/(r^2(2r-1))<0\),
so \(g(r)/(r-1)\) is decreasing. At \(r_0=9/5\) the
rational number
\[
 A=C(1,1/r_0)^{10}=\frac{13^{13}}{18^8\,5^5}
\]
satisfies \(A^5>(31/8)^8\) by integer arithmetic, and
\(\csc(\pi/12)=\sqrt6+\sqrt2<31/8\). Therefore
\(5g(r_0)>(r_0-1)\log\csc(\pi/12)\), and the same
inequality holds for all \(1<r\le r_0\) by concavity.
Together with \eqref{eq:csc5} this proves the degree-five
claim on this interval.

For \(r\ge r_0\), Cauchy--Schwarz in the Gamma integral
gives \(\Gamma(x+\tfrac12)\le\sqrt{x}\,\Gamma(x)\).
Thus \eqref{eq:gamma-bound} gives
\[
 C(5,1/r)\le B_5(r)\le
 \frac{r(r+1)\sqrt{r+2}}{\Gamma(7/2)}=:W_5(r).
\]
Writing \(f(u)=(1/u-1)\log(1-u)\), the log ratio
\(C(1,1/r)^5/W_5(r)\) is a constant plus
\(\tfrac52 f(u)-\log(1+2u)-\tfrac12\log(1+4u)\).
Its derivative is increasing in \(u\), since
\(f'(u)=\sum_{j\ge2}u^{j-2}/j\) is increasing
and the reciprocal terms in the derivative decrease in
magnitude. It is negative at \(u=5/18\): use
\(f'(u)\le1/[2(1-u)]\) to check
\(\tfrac52 f'(u)<2/(1+2u)+2/(1+4u)\) there.
At \(r_0=9/5\) the ratio is greater than one. Squaring both sides
reduces this to
\(A>r_0^2(r_0+1)^2(r_0+2)/\Gamma(7/2)^2\), which follows
using \(\Gamma(7/2)^2=225\pi/64\) and \(\pi>157/50\)
by rational arithmetic. Since \(u\le5/18\) in this range,
the ratio stays greater than one. This proves the
degree-five claim.

For the tail, the function
\(x\mapsto\log\Gamma(r+x)-\log\Gamma(1+x)-\log\Gamma(r)\)
is concave for \(x\ge0\), since the trigamma function decreases.
It vanishes at zero, so \(B_k(r)^{1/k}\) is nonincreasing
in positive \(k\). It suffices to compare \(B_6(r)\) with
\(C(1,1/r)^6\). The ratio is
\[
 \frac{C(1,1/r)^6}{B_6(r)}
   =\frac{48(1-u)^{3/u-3}}{(1+2u)(1+4u)}.
\]
Its logarithmic derivative in \(u\) is increasing, again by
the series for \(f'\). At \(u=1/2\) that derivative is
\(12\log2-25/3<0\). For the strict rational estimate one
can use
\[
 \log2
 \le 2\left(\frac13+\frac1{81}
               +\frac{1}{5\cdot243(1-1/9)}\right)<\frac{25}{36},
\]
obtained by summing the tail of the series for
\(2\,\operatorname{arctanh}(1/3)\).
The ratio equals one at \(u=1/2\), and is therefore strictly
greater for \(u<1/2\). This proves the whole tail by
\eqref{eq:gamma-bound}.
\end{proof}

\begin{theorem}\label{thm:dirichlet-main}
For every \(0<p<1\),
\[
 C(k,p)\le C(1,p)^k\qquad(k\ge1),
\]
with equality exactly at \(k=1,2\).
Consequently the square-free maximal-order conjecture of
Bondarenko--Brevig--Saksman--Seip holds in the original form:
\[
 \limsup_{N\to\infty}\frac{\log D(N,p)}{\log N/\log\log N}
   =\log C(1,p).
\]
Also
\[
 \log\max_{1\le N\le X}D(N,p)
    \sim \log C(1,p)\frac{\log X}{\log\log X}.
\]
\end{theorem}

\begin{proof}
Combine the degree-two identity, Lemmas~\ref{lem:cubic-domination},
\ref{lem:quartic-domination} and \ref{lem:tail-domination}.
Proposition~\ref{prop:extremal-prime-power} gives the limsup.
For the running maximum, the same upper bound is uniform for
all sufficiently large \(N\le X\) since
\(\log N/\log\log N\) is eventually increasing. The largest
primorial not exceeding \(X\) gives the lower bound. Indeed
the prime number theorem implies both that its logarithm is
asymptotic to \(\log X\) and that the number of its prime
factors is asymptotic to \(\log X/\log\log X\).
\end{proof}

\subsection{The first difference from the square-free maximum}

Let \(M_p(X)=\max_{N\le X}D(N,p)\) and
\(M_p^{\mathrm{sf}}(X)=\max_{N\le X,\ N\ {\rm square\!-\!free}}
D(N,p)\). Put \(L=\log X\), \(\ell=\log L\), and
\(B=\log C(1,p)>0\).

\begin{theorem}\label{thm:second-order-dir}
For every fixed \(0<p<1\) we have
\begin{align*}
 \log M_p(X)
   &=B\,\frac{L}{\ell}
        +B(1+\log2)\,\frac{L}{\ell^2}
        +o_p(L/\ell^2),\\
 \log M_p^{\mathrm{sf}}(X)
   &=B\,\frac{L}{\ell}
        +B\,\frac{L}{\ell^2}
        +o_p(L/\ell^2).
\end{align*}
In particular the unrestricted maximum is larger on the
second scale by
\((B\log2+o_p(1))L/\ell^2\) in its logarithm.
\end{theorem}

\begin{proof}
Write \(\vartheta(t)=\sum_{q\le t}\log q\), with sums over primes,
and let \(t\) be the largest prime with \(2\vartheta(t)\le L\).
The square of the associated primorial is at most \(X\).
Since \(C(2,p)=C(1,p)^2\), it gives
\(\log M_p(X)\ge2B\pi(t)\).

The strict inequalities in Theorem~\ref{thm:dirichlet-main}
and the classical bound \(\log C(k,p)=O_p(\log(k+1))\)
give a fixed gap
\[
 \Delta=B-\sup_{k\ge3}\frac{\log C(k,p)}{k}>0.
\]
Set \(\lambda=B/\log t\) and \(\eta=\Delta/(2B)>0\).
For \(N=\prod_q q^{k_q}\le X\), the product formula yields
\[
 \log D(N,p)\le \lambda L+
 \sum_{q}\max_{k\ge0}\{\log C(k,p)-\lambda k\log q\},
 \qquad C(0,p)=1.
\]
The terms vanish for \(q\ge t\). For \(q\ge t^{1-\eta}\), each
\(k\ge3\) gives a strictly negative quantity, by the gap.
The maximum over \(k=0,1,2\) is \(2(B-\lambda\log q)_+\),
using the exact degree-two identity. For the smaller primes,
\(\log C(k,p)\le B+\tau\log k\) for \(k\ge1\). The maximum of
\(B+\tau\log k-\lambda k\log2\) over real \(k\ge1\) is
\(O_p(1+\log(1/\lambda))\). Thus their combined contribution
is \(O_p(t^{1-\eta}\log\log t)\).
The upper bound is consequently
\[
 \lambda L+2B\pi(t)-2\lambda\vartheta(t)
             + O_p(t^{1-\eta}\log\log t).
\]
The next prime is at most \(2t\) for large \(t\), so
\(0\le L-2\vartheta(t)=O(\log t)\).
We obtain \(\log M_p(X)=2B\pi(t)+O_p(t^{1-\eta}\log\log t)\).
The quantitative prime number theorem gives
\(t=L/2+o(L/\ell)\) and
\(\pi(t)=t/\log t+t/(\log t)^2+o(t/(\log t)^2)\).
This proves the first expansion. For square-free \(N\),
\(D(N,p)=e^{B\omega(N)}\). Given \(\omega(N)\), the smallest
possible \(N\) is the corresponding primorial. Thus the
same argument with the largest prime \(t'\) for which
\(\vartheta(t')\le L\) gives the second expansion directly.
\end{proof}

\section{Scope and prior work}

The classical best constant in the \(n^\tau\) bound, the
contractive Bergman embedding, and the exact multiplicative
identity for Dirichlet coefficient norms are established
results, used here in their original normalizations. Our
prime-power comparison closes the question following (34) in
\cite{BBSS}; it makes essential use of the distinguished
third and fourth coefficient formulas of \cite{Low}.
The disk high-degree limit is given by a Fourier extremum
and a two-sided finite-degree enclosure. The latter can
be used to compute it without an assumed closed evaluation
of the Fourier extremum. The literature comparison covers
the references below and targeted searches for the stated
maximal-order question and disk coefficient limit.

\small
\begin{thebibliography}{9}
\setlength{\itemsep}{2pt}
\bibitem{BBSS}
A.~Bondarenko, O.~F.~Brevig, E.~Saksman, and K.~Seip,
Linear space properties of \(H^p\) spaces of Dirichlet series,
\emph{Trans. Amer. Math. Soc.} \textbf{372} (2019), 6677--6702.
\href{https://doi.org/10.1090/tran/7898}{DOI: 10.1090/tran/7898}.

\bibitem{BrevigSaksman}
O.~F.~Brevig and E.~Saksman, Coefficient estimates for \(H^p\)
spaces with \(0<p<1\), \emph{Proc. Amer. Math. Soc.} \textbf{148}
(2020), 3911--3924.
\href{https://doi.org/10.1090/proc/14995}{DOI: 10.1090/proc/14995}.

\bibitem{Kulikov}
A.~Kulikov, Functionals with extrema at reproducing kernels,
\emph{Geom. Funct. Anal.} \textbf{32} (2022), 938--949.
\href{https://arxiv.org/abs/2203.12349}{arXiv:2203.12349}.

\bibitem{Low}
\emph{Sharp third and fourth Taylor coefficients in Hardy
spaces below \(H^1\)}, preceding manuscript, November 2026.

\bibitem{General}
\emph{Sharp Hardy coefficient norms in arbitrary degree},
preceding manuscript, November 2026.
\end{thebibliography}

\appendix
\section{Exact polynomial checks}

The following code verifies the elimination and the starting
sign in Lemma~\ref{lem:quartic-domination}, and the
degree-three and degree-five comparisons. All arithmetic is exact.
\begin{verbatim}
import sympy as S

a,v,t = S.symbols('a v t')
d = a-1
D = a-2+(2*a-1)*v
A2 = a*d/2
A3 = A2*(a-2)/3
A4 = A3*(a-3)/4
k = d*v/(2*(1-v))
u = 6*(1-v)**2*(1+v)/(v*d*D)
S0 = (1+1/a)**4
U = 1+u+(k*u)**2+v*v*u+(v/a)**2
R = v/a*(v+2*A2*v*u+(A2*k*k+3*A3*k+A4)*u*u)-1
N = S.cancel(2*a*v*d*D**2*R)
G = S.cancel(a**3*v*d*D**2*(((a+1)*a*v-2)*U-d*a*v*S0))
assert S.denom(N)==1 and S.denom(G)==1
num = S.resultant(N,G,v)
den = 59049*a**7*(a-2)*d**18*(a+1)**12*(2*a-1)
T = S.Poly(S.cancel(num/den),a)
assert T.degree()==23
assert all(c>0 for c in S.Poly(T.as_expr().subs(a,1+t),t).all_coeffs())
N2 = S.Poly(N.subs(a,2),v)
G2 = S.Poly(G.subs(a,2),v)
assert S.gcd(N2,G2).monic().as_expr()==v
lo,hi = S.Rational(60969,10**5), S.Rational(60970,10**5)
assert N2.eval(lo)>0 and N2.eval(hi)<0
assert sum(c*(lo if c>0 else hi)**j[0] for j,c in G2.terms())>0

r = S.Rational(9,5)
A = 13**13/(S.Integer(18)**8*5**5)
assert A**5 > S.Rational(31,8)**8
assert A > r**2*(r+1)**2*(r+2)/(S.Rational(225,64)*S.Rational(157,50))

B,u,w = S.symbols('B u w')
P = (a-2)*u**3/6+a*(5-a)*u**2/4+a*(1-2*a)*u/2-a*a
assert S.factor(P.subs(u,2*a))==a**3*(a+1)/3
assert S.cancel(P.subs(u,10*a/7)+a*a*(25*a*a-155*a+294)/1029)==0
w = 1+(a-1)*v
U3 = 1+1/(a*a*w)+4*v*v*(w+1)/w**2
S3 = (1+1/a)**3
T3 = 1-(1-v)/a+(1-v)**2/(2*a*a)
J = S.cancel(2*a*a*(a+1)**3*w*w*(U3/S3-T3))
J = J.subs({a:1+B,v:(5+7*t)/(7*(1+t))})
P3 = S.Poly(S.cancel(7**4*(1+t)**4*J),B,t)
assert len(P3.coeffs())==35 and all(c>0 for c in P3.coeffs())
\end{verbatim}

\end{document}
